フィジカル AI（実世界で身体を持って動く AI）のハードウェアは高価で、試行錯誤のコストが大きい。
本報告書は、『機動戦士ガンダム』の SF 考証から生まれた AMBAC（Active Mass Balance Auto Control：手足の質量移動の反作用による姿勢制御）\cite{R1,R42} を出発点に、
安価なコンピュータシミュレーションで、ロボットの形状が姿勢制御・安定性・移動能力に与える意味をどこまで明らかにできるかを検討する。
中心となる問いは「二足歩行は本当にバランスが良いのか、なぜ人は二本脚なのか」である。

剛体力学・幾何力学・制御理論・機械学習を用いた 10 本のシミュレーション（Sim01〜Sim10）により、次のことを示した。
(1) AMBAC は角運動量保存の下で物理的に成立し、手足のループ運動による幾何学的位相として定量的に記述できる。リアクションホイールはその特殊な場合である。
(2) 二足歩行は静的にはバランスが悪いが、動かせる重い上体による能動的なバランス（地上版の AMBAC）と、人間の大きさでの段差への強さを持つ。
形状の優位は、宇宙・平地・段差、人間大・巨大機といった条件で入れ替わり、大型化では二乗三乗則が繰り返し効く。
(3) 人の肩のコッドマンのパラドックスは AMBAC と同じ幾何学的位相であり、腕組みのような人間らしい姿勢には、ひねりを含む 3 自由度の肩と肩甲骨の動きが必要になる。
(4) 学習は AMBAC を発見できるが、物理の構造を方策に組み込むことが成否を分ける。

以上から、用途や目的に応じて形状を選ぶ必要があること、フィジカル AI に必要なスペックとトレードオフを整理した。
最後に、組み合わせた形状について校正値を定義し、要求した動きや性能の実現性を検証する第 2 段階の構想を示す。

\medskip\noindent
\textbf{ソースコード・データ}：\url{\repo}
